\documentclass[a4paper,12pt]{article}
\usepackage{amsmath, amssymb}
\usepackage{xcolor}
\usepackage{tcolorbox}
\usepackage{multicol}
\usepackage{geometry}
\hyphenpenalty=100000

\geometry{left=1.5cm, right=1.5cm, top=2cm, bottom=2cm}

\tcbuselibrary{theorems}

\definecolor{PurpleEx}{HTML}{5d478b}
\definecolor{GrayEx}{HTML}{f2f2f2}
\definecolor{GrayBG}{HTML}{d9d9d9}
\definecolor{BlackSav}{HTML}{000000}

\title{
\centering
\tcbset{colframe=BlackSav,colback=GrayBG}
\tcbox[tikznode]{
Chapitre 3 \\ Divisibilité - Exercices
}
}
\usepackage{tikz}
\author{}
\date{}

\newtcbtheorem
  []% init options
  {exercice}% name
  {Exercice}% title
  {%
    colback=GrayEx,
    colframe=PurpleEx,
    fonttitle=\bfseries,
  }% options
  {ex}% prefix

\begin{document}

\maketitle

\begin{multicols}{2}
% Exercice 1
\begin{exercice}[title=Exercice 1 $\quad \star \quad$ { [Calculer] }]{}{}
    Déterminer les diviseurs positifs des nombres suivants :
    15 ; 27 ; 23 ; 51 ; 11132.
    \end{exercice}
    
    % Exercice 2
    \begin{exercice}[title=Exercice 2 $\quad \star \quad$ { [Calculer] }]{}{}
    Déterminer les entiers naturels $n$ tels que $7$ divise $n + 3$.
    \end{exercice}
    
    % Exercice 3
    \begin{exercice}[title=Exercice 3 $\quad \star \quad$ { [Calculer] }]{}{}
    Déterminer les entiers naturels $n$ tels que $11$ divise $2n + 5$.
    \end{exercice}
    
    % Exercice 4
    \begin{exercice}[title=Exercice 4 $\quad \star \quad$ { [Calculer] }]{}{}
    Déterminer les entiers naturels $n$ tels que $2$ divise $4n - 3$.
    \end{exercice}
    
    % Exercice 5
    \begin{exercice}[title=Exercice 5 $\quad \star \star \quad$ { [Calculer, Raisonner] }]{}{}
    Montrer que pour tout $n \in \mathbb{N}$, $17^n - 2^n$ est divisible par $5$.
    \end{exercice}
    
    % Exercice 6
    \begin{exercice}[title=Exercice 6 $\quad \star \star \quad$ { [Calculer, Raisonner] }]{}{}
    Déterminer l’ensemble des entiers naturels $a$ tels que $a^2 - 24$ soit le carré d’un entier naturel.
    \end{exercice}
    
    % Exercice 7
    \begin{exercice}[title=Exercice 7 $\quad \star \star \quad$ { [Chercher, Représenter] }]{}{}
    Soit $f$ la fonction définie sur $\mathbb{R} \setminus \{1\}$ par $f(x) = \frac{12}{x - 1}$. Déterminer l’ensemble des points de la courbe représentative de $f$ qui sont à coordonnées entières.
    \end{exercice}
    
    % Exercice 8
    \begin{exercice}[title=Exercice 8 $\quad \star \quad$ { [Calculer] }]{}{}
    Déterminer l’ensemble des entiers relatifs $n$ tels que $2n + 5$ divise $7n + 3$.
    \end{exercice}
    
    % Exercice 9
    \begin{exercice}[title=Exercice 9 $\quad \star \quad$ { [Calculer] }]{}{}
    Déterminer l’ensemble des entiers relatifs $n$ tels que $n + 7$ soit divisible par $n$.
    \end{exercice}
    
    % Exercice 10
    \begin{exercice}[title=Exercice 10 $\quad \star \quad$ { [Calculer] }]{}{}
    Déterminer l’ensemble des entiers relatifs $n$ tels que $n + 3$ divise $2n + 5$.
    \end{exercice}
    
    % Exercice 11
    \begin{exercice}[title=Exercice 11 $\quad \star \quad$ { [Calculer] }]{}{}
    Déterminer l’ensemble des entiers relatifs $n$ tels que $4n + 3$ soit divisible par $3n + 4$.
    \end{exercice}
    
    % Exercice 12
    \begin{exercice}[title=Exercice 12 $\quad \star \quad$ { [Calculer, Chercher] }]{}{}
    Déterminer tous les entiers naturels $a$ et $b$ tels que $a^2 - 9b^2 = 24$.
    \end{exercice}
    
    % Exercice 13
    \begin{exercice}[title=Exercice 13 $\quad \star \star \quad$ { [Calculer, Chercher] }]{}{}
    Déterminer tous les entiers naturels $x$ et $y$ tels que $x + y = xy$.
    \end{exercice}
    
    % Exercice 14
    \begin{exercice}[title=Exercice 14 $\quad \star \quad$ { [Calculer] }]{}{}
    Dans chaque cas, déterminer le quotient et le reste de la division de $m$ par $n$.
    \begin{enumerate}
        \item $m = 103$ et $n = 7$
        \item $m = 13$ et $n = 17$
        \item $m = -60$ et $n = 11$
        \item $m = -33$ et $n = -4$
        \item $m = -33$ et $n = 44$
    \end{enumerate}
    \end{exercice}

% Exercice 15
\begin{exercice}[title=Exercice 15 $\quad \star \star \quad$ { [Calculer, Communiquer] }]{}{}
    Dans tout l'exercice, $n$ et $m$ désignent des entiers naturels. Répondre par Vrai ou Faux. Justifier.
    
    \begin{enumerate}
        \item Si le reste de la division euclidienne de $n$ par $7$ est $4$ et le reste de la division euclidienne de $m$ par $7$ est $2$, alors le reste de la division euclidienne de $n + m$ par $7$ est $6$.
        \item Si le reste de la division euclidienne de $n$ par $7$ est $4$ et le reste de la division euclidienne de $m$ par $7$ est $2$, alors le reste de la division euclidienne de $n \times m$ par $7$ est $8$.
        \item Si le reste de la division euclidienne de $n$ par $7$ est $5$ et le reste de la division euclidienne de $m$ par $10$ est $5$, alors le reste de la division euclidienne de $n \times m$ par $17$ est $5$.
        \item Si le reste de la division euclidienne de $n$ par $7$ est $0$ et le reste de la division euclidienne de $m$ par $5$ est $0$, alors le reste de la division euclidienne de $n + m$ par $12$ est $0$.
        \item Si le reste de la division euclidienne de $n$ par $12$ est $1$ et le reste de la division euclidienne de $m$ par $4$ est $2$, alors le reste de la division euclidienne de $n + m$ par $4$ est $3$.
        \item Si le reste de la division euclidienne de $n$ par $15$ est $3$ et le reste de la division euclidienne de $m$ par $5$ est $2$, alors le reste de la division euclidienne de $n + m$ par $15$ est $5$.
    \end{enumerate}
    \end{exercice}
    
    
    % Exercice 16
    \begin{exercice}[title=Exercice 16 $\quad \star \star \quad$ { [Chercher, Calculer] }]{}{}
    Déterminer le reste de la division de $335^{112}$ par $7$.
    \end{exercice}
    
    % Exercice 17
    \begin{exercice}[title=Exercice 17 $\quad \star \quad$ { [Calculer, Raisonner] }]{}{}
    Montrer que pour tout entier naturel $n$, l’entier $n(n+1)(2n+1)$ est divisible par $6$.
    \end{exercice}
    
    % Exercice 18
    \begin{exercice}[title=Exercice 18 $\quad \star \quad$ { [Calculer, Raisonner] }]{}{}
    Montrer que pour tout entier naturel $n$, l’entier $n(n+7)(3n+1)$ est divisible par $3$.
    \end{exercice}
    
    % Exercice 19
    \begin{exercice}[title=Exercice 19 $\quad \star \quad$ { [Calculer, Raisonner] }]{}{}
    Déterminer l’ensemble des entiers naturels $n$ tels que $5n^2 + 3$ soit divisible par $4$.
    \end{exercice}
    
    % Exercice 20
    \begin{exercice}[title=Exercice 20 $\quad \star \star \quad$ { [Chercher, Raisonner] }]{}{}
    Soit $p$ un nombre premier supérieur ou égal à $5$. Montrer que $p^2 - 1$ est divisible par $24$.
    \end{exercice}
    
    % Exercice 21
    \begin{exercice}[title=Exercice 21 $\quad \star \star \star \quad$ { [Chercher, Raisonner] }]{}{}
    Soit $(u_n)_{n \in \mathbb{N}}$ la suite définie par $u_n = 7n - 5$. Déterminer le nombre de carrés parfaits parmi les $2022$ premiers termes de la suite.
    \end{exercice}

    % Exercice 22
\begin{exercice}[title=Exercice 22 $\quad \star \star \star \quad$ { [Chercher, Raisonner] }]{}{}
    Soit $p$ un nombre premier impair. Montrer qu’il existe $n \in \mathbb{N}^*$ tel que $p$ divise $2^n - 1$.
    \end{exercice}
    
    % Exercice 23
    \begin{exercice}[title=Exercice 23 $\quad \star \star \quad$ { [Chercher, Raisonner] }]{}{}
    Soit $n$ un carré parfait dont le chiffre des unités est $5$. Déterminer son chiffre des dizaines. La condition établie sur les chiffres des unités est-elle une condition nécessaire et suffisante ?
    \end{exercice}
    
    % Exercice 24
    \begin{exercice}[title=Exercice 24 $\quad \star \quad$ { [Calculer] }]{}{}
    Dans chaque cas, déterminer si $a$ est inversible modulo $n$ et déterminer l’inverse le cas échéant.
    \begin{enumerate}
        \item $a = 3$ et $n = 7$
        \item $a = 3$ et $n = 10$
        \item $a = 4$ et $n = 12$
        \item $a = 6$ et $n = 9$
        \item $a = 9$ et $n = 16$
    \end{enumerate}
    \end{exercice}
    
    % Exercice 25
    \begin{exercice}[title=Exercice 25 $\quad \star \quad$ { [Calculer] }]{}{}
    Résoudre dans $\mathbb{Z}$ les équations suivantes :
    \begin{enumerate}
        \item $3x \equiv 1 \; [11]$
        \item $3x \equiv 4 \; [12]$
        \item $5x \equiv 7 \; [14]$
        \item $5x \equiv 6 \; [8]$
    \end{enumerate}
    \end{exercice}
    
    % Exercice 26
    \begin{exercice}[title=Exercice 26 $\quad \star \star \quad$ { [Communiquer, Raisonner] }]{}{}
    \begin{enumerate}
        \item Énoncer la règle de divisibilité par $5$ puis la démontrer.
        \item Énoncer la règle de divisibilité par $3$ puis la démontrer.
    \end{enumerate}
    \end{exercice}
    
    % Exercice 27
    \begin{exercice}[title=Exercice 27 $\quad \star \star \star \quad$ { [Raisonner, Communiquer, Chercher] }]{}{}
    La propriété de décomposition d’un entier naturel $n$ en base $10$ s’énonce ainsi :
    Pour tout $n \in \mathbb{N}$, il existe un unique $r \in \mathbb{N}$ et des entiers uniques $a_0, a_1, \ldots, a_r$ tels que
    \[
    0 \leq a_i \leq 9, \quad a_r \neq 0
    \]
    et 
    \[
    n = \sum_{k=0}^r a_k 10^k.
    \]
    \begin{enumerate}
        \item Énoncer de manière similaire la propriété de décomposition d’un entier naturel $n$ en base $2$.
        \item Démontrer la propriété énoncée à la question 1.
        \item Décomposer l’entier $113$ en base $2$.
    \end{enumerate}
    \end{exercice}
    
    % Exercice 28
\begin{exercice}[title=Exercice 28 $\quad \star \star \quad$ { [Communiquer, Calculer] }]{}{}
    Les entiers $1 ; 11 ; 111 ; \dots$ sont appelés des reps-units. On note $N_p$ le rep-unit comprenant $p$ fois le chiffre $1$, c’est-à-dire 
    \[
    N_p = \sum_{k=0}^{p-1} 10^k.
    \]
    \begin{enumerate}
        \item 
        \begin{enumerate}
            \item Montrer que $N_p = \frac{10^p - 1}{9}$.
            \item En déduire que, pour tout $p \geq 1$, $9$ divise $10^p - 1$.
        \end{enumerate}
        \item 
        \begin{enumerate}
            \item Déterminer le reste de la division de $10^p$ par $7$ suivant les valeurs de $p$.
            \item En déduire les valeurs de $p$ pour lesquelles $N_p$ est divisible par $7$.
        \end{enumerate}
        \item Montrer que si $p$ est pair, $N_p$ est divisible par $11$.
    \end{enumerate}
    \end{exercice}
    
    % Exercice 29
    \begin{exercice}[title=Exercice 29 $\quad \star \star \quad$ { [Modéliser, Chercher] }]{}{}
    En informatique, un octet est composé de huit bits ($0$ ou $1$). Plus précisément, il est composé de sept bits contenant l’information à transmettre auquel on ajoute un bit de parité de la façon suivante : on calcule la somme $s$ des sept bits constituant l’octet et on prend $m$ tel que $s + m \equiv 0 \; [2]$.
    \begin{enumerate}
        \item Déterminer le bit de parité de $1011101$.
        \item Démontrer que si un seul des sept bits contenant l’information est modifié, le bit de parité détecte l’erreur.
        \item Dans quel cas les erreurs de transmission seront-elles détectées ? Justifier.
    \end{enumerate}
    \end{exercice}
    
    % Exercice 30
    \begin{exercice}[title=Exercice 30 $\quad \star \star \quad$ { [Calculer, Modéliser] }]{}{}
    On souhaite chiffrer un message. On procède de la manière suivante :
    \begin{itemize}
        \item Chaque lettre du texte à chiffrer est associée à un entier $x$ entre $0$ et $25$ selon la correspondance suivante :
        \[
        \begin{array}{ccccccccccccc}
        A & B & C & D & E & F & G & H & I & J \\
        0 & 1 & 2 & 3 & 4 & 5 & 6 & 7 & 8 & 9 \\
        \end{array}
        \]
        \[
        \begin{array}{cccccccccccc}
         K & L & M & N & O & P & Q & R & S \\
         10 & 11 & 12 & 13 & 14 & 15 & 16 & 17 & 18 \\
        \end{array}
        \]
        \[
        \begin{array}{ccccccc}
        T & U & V & W & X & Y & Z \\
        19 & 20 & 21 & 22 & 23 & 24 & 25 \\
        \end{array}
        \]
        \item À chaque entier $x$, on associe l’entier $y$ tel que $y = 21x + 2 \; [26]$ avec $0 \leq y < 25$.
        \item On remplace la lettre du message par la lettre correspondant au nombre $y$.
    \end{itemize}
    \begin{enumerate}
        \item Montrer que le mot MATHS est chiffré par le mot UCLTQ.
        \item L’objectif de cette question est de déchiffrer le mot suivant :
        \[
        \text{LGVOPY}
        \]
        \begin{enumerate}
            \item Montrer que si $x$ est un entier et $y$ l’entier qui lui est associé par chiffrement, alors $x = 5y - 10 \; [26]$.
            \item En déduire le message clair correspondant à LGVOPY.
        \end{enumerate}
    \end{enumerate}
    \end{exercice}
    
    % Exercice 31
    \begin{exercice}[title=Exercice 31 $\quad \star \star \star \quad$ { [Modéliser, Chercher] }]{}{}
    Le numéro INSEE d’une personne est inscrit sur sa carte vitale. Il est composé de $13$ chiffres correspondant à :
    \begin{itemize}
        \item le sexe ($1$ pour un homme et $2$ pour une femme) ;
        \item l’année de naissance (deux chiffres) ;
        \item le mois de naissance (deux chiffres) ;
        \item le lieu de naissance (cinq chiffres correspondant au département et à la commune) ;
        \item le numéro d’ordre d’inscription des naissances dans la commune (trois chiffres).
    \end{itemize}
    Une clé de contrôle à deux chiffres complète le numéro INSEE. On calcule la clé avec la méthode suivante : on calcule le reste $r$ de la division du nombre à $13$ chiffres de la carte par $97$ et la clé est alors égale à $97 - r$.
    \begin{enumerate}
        \item Déterminer la clé de contrôle correspondant au numéro INSEE suivant : $254025002500522$.
        \item On suppose que lors d’une saisie d’un code INSEE, une erreur est commise.
        \begin{enumerate}
            \item Montrer que si un seul chiffre est erroné, l’erreur est détectée par la clé.
            \item Montrer que si l’on intervertit les deux premiers chiffres, l’erreur est détectée par la clé.
        \end{enumerate}
    \end{enumerate}
    \end{exercice}    

    % Exercice 32
    \begin{exercice}[title=Exercice 32 $\quad \star \star \star \quad$ { [Modéliser, Chercher, Communiquer] }]{}{}
    Faire une recherche sur le système EAN-13 utile pour les codes-barres et montrer que le système permet de détecter certaines erreurs.
    \end{exercice}

\end{multicols}

\end{document}
